\documentclass[12pt]{report}
% Definiciones basicas de LaTeX para la memoria
%
% No modificar las lineas a continuacion a menos
% de estar muy seguro de lo que se esta haciendo!!
%
% Jaime Cisternas, julio 2004, agosto 2010

% <jcisternas@uandes.cl> 

%%%%%%%%%%%%%%%%%%%%%%%%%%%%%%%%%%%%%%%%%%%%%%%%%%%%%%%%%%%%%%%%%%%%%
% Libreria computin
%\usepackage{algorithm}
%\usepackage{algorithmic}
\usepackage[left=3cm,top=2.5cm,right=3cm,bottom=2.5cm]{geometry}

% Encoding

%\usepackage[utf8]{inputenc} % unnecessary when compiling through xelatex
%\usepackage[T1]{fontenc}

% Carga librerias utiles con simbolos y estructuras
\usepackage{amsfonts}
\usepackage{amssymb}
\usepackage{amsmath}   
\usepackage{amsthm}
\usepackage{latexsym}
\usepackage{cancel}
\usepackage{physics}
\usepackage{bbm}
\usepackage{listings}
\usepackage{pgfplots}
\usepackage{eurosym}

% Logo Julia

% \usepackage{core/julialogo}

% Cool code highlight
% https://es.overleaf.com/learn/latex/Code_Highlighting_with_minted
% No need of outputdir out of overleaf
\usepackage[outputdir=docs/DocumentoMemoria]{minted}

% Librerías adicionales para diagramas tikz

\usepackage{tikz}
\usepackage{mathdots}
\usepackage{yhmath}
\usepackage{color}
\usepackage{xcolor}
\usepackage{siunitx}
\usepackage{multirow}
\usepackage{gensymb}
\usepackage{tabularx}
\usepackage{extarrows}
\usepackage{booktabs}
\usetikzlibrary{fadings}
\usetikzlibrary{patterns}
\usetikzlibrary{shadows.blur}
\usetikzlibrary{shapes}
\usepackage{comment} 

% Libreria para ajustar espaciado entre lineas
\usepackage{setspace}

% Libreria con estilo APA para bibliografia
%\usepackage{apacite}
%\bibliographystyle{apacite}
\usepackage{core/newapa}

% Si usted está utilizando un teclado en español, puede ser útil
% usar los paquetes a continuacion. De otra manera tendrá que generar las
% tildes y la letra ñ de forma especial.
% Estos caracteres no pueden ser utilizados en etiquetas ni tampoco en modo matemático
\usepackage[spanish, activeacute]{babel}
\unaccentedoperators


% Declaraciones especificas de pdfLaTeX
\usepackage[%pdftex, % commented to compile with Xelatex
        colorlinks=false,         % true or false (for final version)
        urlcolor=rltblue,         % \href{...}{...} external (URL)
	    anchorcolor=rltbrightblue,
        filecolor=weben,          % \href*{...} local file
        linkcolor=webred,         % \ref{...} and \pageref{...}
        menucolor=webdarkblue,
        citecolor=webgreen,
        pdftitle={},              % INSERT YOUR TITLE HERE
        pdfauthor={},             % INSERT YOUR NAME HERE
        pdfsubject={},
        pdfkeywords={},
        %pdfpagemode=None,
        bookmarksopen=true,
	plainpages=false]{hyperref}
%\usepackage[pdftex]{graphicx}

\usepackage{graphicx}

\usepackage{fontspec}
 
\setmainfont{Times New Roman}



%\pdfcompresslevel=9 %commented to use Xelatex
%\pdfadjustspacing=1 %commented to use xelatex
% Definicion de colores
\usepackage{color}
\definecolor{rltbrightred}{rgb}{1,0,0}
\definecolor{rltred}{rgb}{0.75,0,0}
\definecolor{rltdarkred}{rgb}{0.5,0,0}
\definecolor{rltbrightgreen}{rgb}{0,0.75,0}
\definecolor{rltgreen}{rgb}{0,0.5,0}
\definecolor{rltdarkgreen}{rgb}{0,0,0.25}
\definecolor{rltbrightblue}{rgb}{0,0,1}
\definecolor{rltblue}{rgb}{0,0,0.75}
\definecolor{rltdarkblue}{rgb}{0,0,0.5}
\definecolor{webred}{rgb}{0.5,.25,0}
\definecolor{webblue}{rgb}{0,0,0.75}
\definecolor{webgreen}{rgb}{0,0.5,0}
\definecolor{webdarkblue}{rgb}{0,0,0.5}
\definecolor{webbrightgreen}{rgb}{0,0.75,0}
% fin de declaraciones especificas pdfLaTeX

% tamanio de pagina
%\setlength{\oddsidemargin}{.5in}
%\setlength{\evensidemargin}{.0in} % en caso de usar opcion 'twoside'
%\setlength{\textwidth}{6in}
%\setlength{\topmargin}{-.5in}
%\setlength{\textheight}{9in}


% Las siguientes definiciones pueden ser usadas en
% memorias mas matematicas
\newtheorem{theorem}{Teorema}[section]
\newtheorem{lemma}[theorem]{Lema}
\newtheorem{corollary}[theorem]{Corolario}
\newtheorem{proposition}[theorem]{Proposición}
\newtheorem{definition}[theorem]{Definición}
\newtheorem{claim}{Afirmación}
\newtheorem{conjecture}[theorem]{Conjetura}
\newtheorem{observation}[theorem]{Observación}
\newtheorem{problem}[theorem]{Problema}

% Definicion de un ambiente para algoritmos en pseudo-codigo.
% Esta definicion puede ser mejorada.
\newtheorem{algorithm}{Algoritmo}[section]
\newcommand{\tab}{\hspace*{0.5 cm}}
% Palablas claves que deben aparecer en negrita
\newcommand{\bfWHILE}{\textbf{while~}}
\newcommand{\bfDO}{\textbf{do~}}
\newcommand{\bfEND}{\textbf{end~}}
\newcommand{\bfIF}{\textbf{if~}}
\newcommand{\bfTHEN}{\textbf{then~}}
\newcommand{\bfELSE}{\textbf{else~}}
\newcommand{\bfFOR}{\textbf{for~}}

% Nombres fijos de Latex pueden ser cambiados al Espaniol
\renewcommand{\contentsname}{Índice General}
\renewcommand{\chaptername}{Capítulo}
\renewcommand{\appendixname}{Anexo}
\renewcommand{\bibname}{Bibliografía}
\renewcommand{\figurename}{Ilustración}
\renewcommand{\tablename}{Tabla}
\renewcommand{\indexname}{Índice}
\renewcommand{\partname}{Parte}
\renewcommand{\listfigurename}{Lista de Ilustraciones}
\renewcommand{\listtablename}{Lista de Tablas}
\AtBeginDocument{\RenewCommandCopy\qty\SI}

% Para las enumeraciones usamos primero a,b,c,... y despues i, ii, iii,...
\renewcommand{\labelenumi}{\alph{enumi})}
\renewcommand{\labelenumii}{\roman{enumii})}
\renewcommand{\labelenumiii}{-}
\renewcommand{\labelenumiv}{-}
\pgfplotsset{compat=1.18}

% Para ser consistentes en el uso de abreviaciones!
%\newcommand{\Chp}{Cap\'{\i}tulo}
%\newcommand{\Chps}{Cap\'{\i}tulos}
%\newcommand{\Sec}{Sec.} % \S
%\newcommand{\Secs}{Secs.} % \S\S
%\newcommand{\SSec}{Subsec.}
%\newcommand{\Fig}{Fig.}
%\newcommand{\Figs}{Figs.}
%\newcommand{\Eqn}{Ecn.}
%\newcommand{\Eqns}{Ecns.}

% Muestra en pantalla los nombres de los archivos usados
\listfiles

% Previene una cierta senial de ``warning: contentsline with no destination''
\newcounter{dummy}



\setlength{\parindent}{0cm}
\usepackage{float}
\usepackage{subfigure}
\usepackage{array}
\newcolumntype{E}{>{$}c<{$}}
\usepackage{longtable}
\setcounter{MaxMatrixCols}{40}
\usepackage{bm}

%%%%%%%%%%%%%%%%%%%%%%%%%%%%%%%%%%%%%%%%%%%%%%%%%%%%%%%%%%
\newcommand{\nombreautor}{Samuel Honorato Rojas}
\newcommand{\mes}{Octubre}
\newcommand{\anio}{2023}
\newcommand{\titulo}{Evaluación de la viabilidad económica de la transición hacia combustibles renovables en Chile: un análisis basado en un modelo de equilibrio con múltiples productores}
\newcommand{\nombreprofuno}{Sebastián Cea Echenique}
\newcommand{\nombreprofdos}{Nombre Profe 2}
\newcommand{\nombreproftres}{Nombre Profesor Invitado}
\newcommand{\codigo}{ING-IN-001/11}

%%%%%%%%%%Aqui se escribe el resumen%%%%%%%%%%%%%%%%%%%%%
\newcommand{\resumen}{
Hacer resumen al terminar la tesis. Máximo 1 página
}

\newcommand{\agradecimientos}{
Muchas gracias a todos}

\newcommand{\dedicatoria}{\it
ok..
}
\renewcommand{\baselinestretch}{1.5}

%%%%%%%%%%%%%%%%%%%%%%%%%%%%%%%%%%%%%%%%%%%%%%%%%%%%%%%%%%

% Por favor coloque todas las definiciones de simbolos a continuacion
% y no en los archivos de capitulos. Esto evitara la  existencia de
% multiples definiciones para las mismas palabras.

\newcommand{\mydef}
	{\stackrel{\mathrm{def}}{=}}
\newcommand{\e}
	{\hbox{\large{e}}}
\newcommand{\mi}
	{\hbox{\large{i}}}
\newcommand{\RR}{\mathbb{R}}


%%%%%%%%%%%%%%%%%%%%%%%%%%%%%%%%%%%%%%%%%%%%%%%%%%%%%%%%%%

\begin{document}
\label{start}

% Genera las paginas de titulo, copyright, resumen,
% agradecimientos, dedicatoria, indices,...

% No es necesario modificar

% Paginas de titulo, derechos de autor y otras mas.

% Jaime Cisternas, julio 2004, agosto 2010

% declaracion especifica de pdfLaTeX
% permite incluir figuras con \includegraphics[...]{figurename}
\DeclareGraphicsExtensions{.jpg,.pdf,.mps,.png}

%%%%%%%%%%%%%%%%%%%%%%%%%%%%%%%%%%%%%%%%%%%%%%%%%%%%%%%%%%%%%%%%

  \pagenumbering{roman}
% \setcounter{page}{1}

% Escribe la pagina de titulo, incluyendo el autor y la facultad

  \thispagestyle{empty}
  \textsc{
  \vspace*{0cm}
  \begin{center}
    \Large
     Universidad de los Andes \\
     Facultad de Ingeniería y Ciencias Aplicadas\\
  \end{center}
     \vspace{1cm}
  \begin{center}
     \includegraphics[angle=0,height=5cm]{Samuel/Memoria Samuel/Imagenes/Logo-UANDES.png}
  \end{center}
     \vspace{1cm}
  \begin{center}
     \Large
     \titulo
  \end{center}
  \vspace{0.5cm}
  \begin{center}
    \Large
    \nombreautor
  \end{center}
  \vspace{0.5 cm}
  \begin{center}
    Memoria para optar al título de \\
    Ingeniero Civil Industrial 
    Mención Gestión de Operaciones\\
    \vspace{1cm}
    Profesor Guía: \nombreprofuno \\
    \vspace{0.5cm}
    \codigo \\
    \vspace{0.5cm}
    Santiago, \mes\ de \anio
  \end{center}
  }

%%%%%%%%%%%%%%%%%%%%%%%%%%%%%%%%%%%%%%%%%%

% Escribe la segunda pagina de titulo para ser firmada por la comision


\thispagestyle{empty}

\begin{center}

\vspace*{2cm}
\parbox{10cm}{
\noindent
Certifico que he leído esta memoria y que en mi opinión
su alcance y calidad son completamente adecuados como para ser considerada
una memoria conducente al título de Ingeniero.
\vspace{1cm}

\hfill
\begin{tabular}{c}
\hspace{8cm} \\
\hline
\nombreprofuno \\
(Profesor Guía)
\end{tabular}

\vspace*{1.5cm}

}

\end{center}

%%%%%%%%%%%%%%%%%%%%%%%%%%%%%%%%%%%%%%%%%%

% Espaciado entre lineas
  \onehalfspacing
  %\doublespacing

% Hace pagina de derechos de autor


  \thispagestyle{empty}
  %\vspace*{0in}
  \begin{center}
    \copyright\ \nombreautor\ \anio \\
    Todos los derechos reservados.
  \end{center}

%%%%%%%%%%%%%%%%%%%%%%%%%%%%%%%%%%%%%%%%%%%%%%%%%%%%%%%%%%%%%%%%%%%%%%

% Genera agradecimientos leyendo definicion de agradecimientos

  \cleardoublepage \addcontentsline{toc}{section}{Agradecimientos}
  \begin{center} \Large \textbf{Agradecimientos} \end{center}


\vspace{2.5mm}




% Genera resumen leyendo definicion de resumen



  \refstepcounter{dummy} \addcontentsline{toc}{section}{Resumen}
  \begin{center} \Large \textbf{Resumen} \end{center}

\vspace{2.5mm}


\vspace{2.5mm}





% Genera dedicatoria

 % \cleardoublepage \vspace*{1.5in}
  %\begin{flushright} \dedicatoria %\end{flushright}

% Produce indice general, listas de figuras y de tablas



  \normalsize
  \pagenumbering{arabic}

%%%%%%%%%%%%%%%%%%%%%%%%%%%%%%%%%%%%%%%%%%%%%%%%%%%%%%%%%%%

\tableofcontents
%%%%%%%%%%%%%%%%Ingreso de capitulos%%%%%%%%%%%%%%%%%%%%%%

% Intro.tex

\chapter{Introducción}
\label{c1} % la etiqueta para referencias




El European Green Deal es una iniciativa de la Unión Europea que busca transformar su economía en una más sostenible y neutral en carbono. El objetivo mas próximo de esta es reducir las emisiones de gases de efecto invernadero en al menos un 55\% para el año 2030 en comparación con los niveles de 1990 \citeB{comision_europea_comunicacion_2019}\\

Para lograr este objetivo, establecieron una serie de medidas y políticas en diferentes sectores, incluyendo la industria del transporte y la producción de energía. En el sector del transporte, se promueve la electrificación de los vehículos y la reducción de las emisiones de los vehículos de combustión interna. Además, se busca fomentar el uso del transporte público y la movilidad sostenible.\\

En el sector de la producción de energía, el European Green Deal busca promover la transición hacia fuentes de energía renovable y la eliminación gradual de los combustibles fósiles. Para lograr este objetivo, se están implementando medidas para reducir la dependencia de los combustibles fósiles y promover el uso de fuentes de energía renovable, cómo la energía eólica y solar, o los distintos combustibles renovables cómo por ejemplo el hidrógeno verde.\\

La Estrategia Europea del Hidrógeno es una de las principales políticas que se desarrollan en el marco del pacto. Se están invirtiendo en tecnologías de producción de hidrógeno verde más eficientes y se están estableciendo incentivos para su uso en diferentes sectores, como el transporte, la producción de energía y las industrias en general. Para esto existe una estrategia que establece los principales campos de actuación para desarrollar el papel del hidrógeno limpio en la reducción de emisiones de la economía de la UE de manera eficiente. La Estrategia establece tres horizontes temporales (2024, 2030 y 2050) para los que determina una sucesión de hitos a alcanzar.\\

En la primera fase (2020-2024), se instalarán al menos 6 GW de electrolizadores en la UE y se producirá hasta 1 millón de toneladas de hidrógeno renovable. En la segunda fase (2025-2030), se espera instalar al menos 40 GW de electrolizadores para 2030 y producir hasta 10 millones de toneladas de hidrógeno renovable en la UE. En la tercera fase (2030-2050), se espera que las tecnologías de hidrógeno renovable alcancen la madurez y se desplieguen a gran escala para llegar a todos los sectores difíciles de descarbonizar.\\

La Estrategia Europea del Hidrógeno establece que el hidrógeno renovable es esencial para respaldar el compromiso de la UE de alcanzar la neutralidad de carbono en 2050 y para respaldar el esfuerzo global para implementar el Acuerdo de París. Además, indica que el ecosistema del hidrógeno en Europa es probable que se desarrolle a través de una trayectoria gradual, a diferentes velocidades en los diferentes sectores y posiblemente en las distintas regiones, requiriendo, por tanto, múltiples soluciones políticas.\\




En este contexto, la transición hacia fuentes de energía renovable se ha convertido en una alternativa cada vez más atractiva y necesaria. En Chile, la producción y consumo de energía son fundamentales para el desarrollo económico y social del país, por lo que es crucial evaluar la viabilidad económica de la transición hacia combustibles renovables.\\

La transición hacia fuentes de energía renovable es la única forma de garantizar un futuro energético sostenible y seguro para Chile. Actualmente, Chile tiene una gran dependencia de los combustibles fósiles, que representan alrededor del 70\% de la matriz energética chilena \citeB{bernal_matriz_nodate}. Además, Chile es vulnerable a los efectos del cambio climático, lo que subraya la importancia de reducir las emisiones de gases de efecto invernadero.\\

Por lo tanto, la transición hacia fuentes de energía renovable es crucial para Chile, no solo para reducir su huella de carbono, sino también para mejorar la seguridad energética y fomentar el desarrollo sostenible. Esta transición es una oportunidad para Chile de liderar la transición energética en América Latina. Sin embargo, esta transición no es sencilla y requiere una planificación cuidadosa, la inversión en tecnologías y la implementación de políticas efectivas.\\

En este sentido, la Corporación de Fomento de la Producción (CORFO) tiene sus principios apuntando en la misma dirección. José Miguel Benavente el Vicepresidente ejecutivo de la CORFO comenta que se ha creado un comité de esta organización. En esta se tienen una nueva estrategia nacional del compromiso de descarbonización, en la que se vincula a los temas industriales del hidrógeno verde en Chile. Este comité lo lidera el Ministerio de Energía.\\

Chile acaba de firmar un acuerdo para que de aquí a 3 años más seamos la mayor producción de hidrógeno verde del mundo. Actualmente la CORFO está atrayendo inversionistas para realizar esta transición a producir hidrógeno verde. Esto lo hace de dos maneras; transformar la industria local para incorporar hidrógeno verde y la otra tiene que ver con la oferta, la parte productiva de electrolizadores.\\

Pretende adelantarse en la carrera de la producción de electrolizadores, comenta que "Podemos ser marginalmente más caros, incluso corrigiendo por los costos de transporte, pero si tenemos una energía que es más barata en términos relativos, energía verde, obviamente hay una compensación entre uno y otro y por lo tanto eso no se transmite al precio final del hidrógeno verde"\citeB{benavente_hidrogeno_2023}.\\

Con esto podemos concluir que aunque el costo de producción del hidrógeno verde en Chile podría ser ligeramente más alto que en otros lugares, esto podría ser compensado por el hecho de que la energía utilizada para producirlo es más barata en términos relativos, ya que proviene de fuentes de energía renovable y verde. Por lo tanto, el precio final del hidrógeno verde no se vería afectado por el costo de producción marginalmente más alto, ya que la energía utilizada para producirlo es más barata y compensa ese costo adicional. Todo esto permitiría al país tener una mayor independencia energética.\\

Además, existen diversos estudios que demuestran el potencial de Chile para la producción de biocombustibles (BC) mediante diferentes métodos. Un ejemplo de esto es un estudio llevado a cabo por la Universidad Técnica Federico Santa María \citeB{ortega_evaluacion_2010}, el cual señala que la opción más viable para producir BC en Chile es el bioetanol a partir del maíz o biodiesel utilizando raps. Sin embargo, es importante mencionar que esta alternativa implicaría un costo superior al de importar directamente los biocombustibles. Pero al igual que lo ya mencionado anteriormente, estaría brindaría al país la ventaja de lograr una mayor diversificación en su matriz energética.\\

 También tenemos el informe \citeB{ortega_evaluacion_2007}, este también concluye que es factible la producción de biodiesel  a partir de aceites vegetales de primer uso y el bioetanol a partir de diferentes cultivos, como la caña de azúcar, la remolacha y el maíz respectivamente.\\


Considerando el enorme potencial que Chile tiene para la producción de biocombustibles y el fuerte respaldo del gobierno, es probable que las diversas industrias se sientan motivadas a realizar una transición hacia los combustibles renovables. Los productores de energía tendrán la opción de generar energía limpia, mientras que la industria del transporte podrá utilizar combustibles limpios, y lo mismo sucederá en todas las demás industrias.\\

\section{Motivación}
Aunque la transición hacia fuentes de energía renovable se presenta como una oportunidad de gran valor para el país, también plantea numerosos desafíos. Los costos de inversión iniciales son altos y las tecnologías aún están en desarrollo. Además, la producción y utilización de biocombustibles y el hidrógeno verde implican cuestiones técnicas y regulatorias que deben ser abordadas.\\

El costo de los combustibles fósiles y sus consecuencias ambientales, unidos a la necesidad de reducir la dependencia energética y de cumplir con los objetivos de descarbonización, hacen que el desarrollo de los combustibles renovables sea imprescindible. Sin embargo, para poder implementar eficazmente esta transición, es necesario llevar a cabo un análisis en profundidad de los factores económicos, técnicos y regulatorios involucrados.\\

Este trabajo tiene como objetivo proporcionar un análisis integral de la viabilidad económica de la transición hacia combustibles renovables en Chile. Con este estudio se pretende proporcionar una herramienta que permita a los encargados de tomar decisiones evaluar las implicaciones económicas de la transición hacia los combustibles renovables, y así desarrollar políticas y estrategias eficaces para fomentar su uso.\\

Estoy convencido de que es necesario trabajar en soluciones sostenibles con el medio ambiente, y esta memoria es mi pequeña contribución en ese sentido. Espero que los hallazgos y conclusiones de este trabajo puedan servir como base para la implementación de políticas y acciones concretas en favor del medio ambiente y el desarrollo sostenible de Chile.\\


\section{Objetivos}
\subsection{Objetivo general}
El objetivo general del presente trabajo consiste en evaluar la viabilidad económica de la transición hacia combustibles renovables en el marco teórico del modelo de \textit{cap and trade} de dos etapas con multiples productores, creado por \citeB{amigo_two_2021}.\\

\subsection{Objetivos específicos}

\begin{enumerate}
\item 
\item 
\item 
\item 
\item 
\item 
\end{enumerate}


\section{Alcances}




\section{Metodología}



\newpage
\section{Estructura del documento}

La estructura del documento se divide en cinco capítulos: introducción, marco teórico, desarrollo metodológico, análisis de resultados y conclusiones.



\input{Samuel/Memoria Samuel/Capitulos/Marco Teórico}
\input{Samuel/Memoria Samuel/Capitulos/Desarrollo Metodológico}
% Ingenieros prefieren la bibliografia antes que los anexos

\cleardoublepage


%\addcontentsline{toc}{Capitulos}{Bibliografía}

\bibliographystyle{newapa}
\bibliography{referencias}

\appendix % seguido de archivos de apendices

\cleardoublepage

% Escribe 'Anexos' en el Indice General
% \addcontentsline{toc}{Capitulos}{Anexos}

% anexo_a.tex

\chapter{Anexos}
\label{}

\section{Códigos Resolución Maere 3.3.1}

\subsection{Código Resolución con Python}\label{codigoPhyton}
\begin{footnotesize}
\begin{lstlisting}
from gurobipy import Model, GRB, quicksum

def modeloRNcompetitivecapacityequilibrium(cant_esce, a, theta, B, tau, c, I):
    m = Model()

    # Decision variables x, Y, Q
    x = m.addVar(lb=0, vtype=GRB.CONTINUOUS, name='x')  # Capacidad
    Y = {}  # Producción
    Q = {}  # Demanda

    # Definición de Producción y Demanda para cada estado de la naturaleza
    for i in range(cant_esce):
        Y[i] = m.addVar(vtype=GRB.INTEGER, name='Y{}'.format(i))
        Q[i] = m.addVar(vtype=GRB.INTEGER, name='Q{}'.format(i))

    # Restricciones
    for i in range(cant_esce):
        # Production equals demand
        m.addConstr(Y[i] == Q[i])
        # Production bounded by installed capacity
        m.addConstr(Y[i] <= x)

    # F.O
    obj_expr = I * 1000 * x + tau * quicksum(theta[i] * (c * Y[i] - a[i] * Q[i] 
    + (B / 2) * (Q[i] * Q[i])) for i in range(cant_esce))
    m.setObjective(obj_expr, GRB.MINIMIZE)

    m.optimize()

    P = []
    IM = []

    for i in range(cant_esce):
        P.append(a[i] - B * Q[i].x)
        IM.append((tau * (P[i] - c) * Y[i].x - I * x.x) / x.x)

    Qprom = sum(Q[i].x for i in range(cant_esce))

    print("Q Promedio =", Qprom / cant_esce)
    print("x =", x.x)
    print("Valor función objetivo =", m.objVal)

    print(cant_esce, a, theta, B, tau, c, I)

    print("\nQ, Y, P and Investment Margin values:")
    for i in range(cant_esce):
        print("Q({}) = {:.1f}    Y({}) = {:.1f}    P({}) = {:.1f}    
        Investment Margin = {:.1f}"
          .format(i + 1, Q[i].x, i + 1, Y[i].x, i + 1, P[i], IM[i]))
\end{lstlisting}
\end{footnotesize}

\subsection{Código Resolución con Julia}\label{codigoJulia}
\begin{footnotesize}
\begin{lstlisting}
using JuMP, Ipopt
A = [300, 350, 400, 450, 500]
theta = [0.2, 0.2, 0.2, 0.2, 0.2]
B = 1
I = 90
C = 60
tau = 8760

model = Model(Ipopt.Optimizer)

# Variables
@variable(model, x >= 0)  # Capacidad
@variable(model, Y[1:5] >= 0) # Producción en escenario w
@variable(model, Q[1:5] >= 0) # Demanda en escenario w

# Restricciones
@constraint(model, demanda[i = 1:5], Y[i] == Q[i])
@constraint(model, capacidad[i = 1:5], Y[i] - x <= 0)

# Función Objetivo
@objective(model, Min, I*1000*x + tau * sum(theta[i] * (C * Y[i] - A[i] * 
Q[i] + B * 0.5 * Q[i]^2) for i = 1:5))
optimize!(model)

# Resultados
println("Q = ", value.(Q))
Qlist = value.(Q)

j = 1
P = []
while j < 6
    Pi = A[j] - B * Qlist[j]
    append!(P, Pi)
    j = j + 1
end
println("P = ", P)
println("x = ", value(x))
\end{lstlisting}
\end{footnotesize}    % archivo de apendice

%\input{apendice_b}    % otro archivo de apendice

\label{end}
\end{document}

%%%%%%%%%%%%%%%%%%%%%%%%%%%%%%%%%%%%%%%%%%%%%%%%%%%%%%%%%%%%